\chapter{Evaluation}
\label{chap:Evaluation}

this chapter evaluates the work against the original aim of the dissertation. because the original python generator was already capable of producing satisfactory RTEMS tests, the question here is not whether the haskell port creates a better test suite, but whether it preserves the existing one while fitting into the same workflow.

the evaluation therefore focuses on two things. the first is output parity between the python and haskell refine command across the seven models listed in the repository: barrier-mgr, chains, event-mgr, msg-mgr, sem-mgr, task-mgr, and proto-sem. the second is workflow compatibility, meaning that the haskell version can be used in the existing pipeline without manual fixes. a second experiment on ease of extension had originally been planned, but was not completed in time and is therefore treated as discussion rather than evidence.

\section{Experiments}

for each of the seven models, test code is generated once with the original python refine command and once with the haskell version. the generated files are then compared using the diff script. any differences are examined and classified as either semantic differences, which change the meaning of the generated test, or non-semantic differences, such as formatting.

the run is counted as successful only if no manual intervention is required and any remaining differences do not affect the meaning of the generated test. this makes the evaluation narrow, but it matches the central risk of the project, which is regression in generated output rather than raw performance.


\section{Results}

this gives a clear success criterion for the dissertation. if the haskell refiner can be substituted into the existing workflow and still generate equivalent tests for the full model set, then the main aim of the project has been met. this does not by itself prove that the new code is easier to extend. that claim is supported more weakly by the implementation itself, particularly the clearer handling of state and the separation between parsing and refinement.

the incomplete extension experiment is therefore a limitation of the evaluation. the project can make a strong claim about behavioural preservation, but only a weaker argued claim about maintainability.



\section{Summary}

the evaluation is based primarily on behavioural preservation. the important question is whether the haskell refiner can replace the python one without changing the generated tests or requiring manual adjustment to the workflow. this is a narrower evaluation than originally planned, but it is also the one most closely tied to the main aim of the project.